\documentclass[../../../include/open-logic-section]{subfiles}

\begin{document}

\olfileid[pt]{sth}{cardinals}{cp}
\olsection{Princípio de Cantor}

Volte o pensamento a \olref[ordinals][vn]{sec}. Estávamos a discutir
conjuntos bem ordenados e sugerimos que seria desejável ter objetos
que funcionem como representantes dos bons ordenamentos. Com isto em mente, introduzimos
os ordinais e depois mostramos em
\olref[ordinals][ordtype]{ordtypesworklikeyouwant} que estes
se comportam como desejamos, isto é:
\[
	\ordtype{A, <} = \ordtype{B, \lessdot}
	\text{ sse } \tuple{A, <} \isomorphic \tuple{B, \lessdot}.
\]
Volte ainda mais atrás, a \olref[sfr][siz][equ]{sec}.
Ali, trabalhando ingenuamente, introduzimos a noção de ``tamanho'' de
um conjunto. Especificamente, dissemos que dois conjuntos são equipotentes,
$\cardeq{A}{B}$, precisamente quando existe !!a{bijection} $f \colon A \to
B$. Esta noção é intrinsecamente {mais simples} do que a de um
bom ordenamento: interessam-nos apenas as !!{bijection}s, e não (como
com os isomorfismos de ordem) se as !!{bijection}s ``preservam alguma
estrutura''.

Tudo isto sugere um pensamento óbvio. Tal como introduzimos
certos objetos, \emph{ordinais}, para calibrar bons ordenamentos, podemos
introduzir certos objetos, \emph{cardinais}, para calibrar tamanho. Esse
é o objetivo deste capítulo.

Antes de dizer o que serão estes cardinais, convém enunciar um
princípio que devem satisfazer. Escrevendo $\card{X}$ para a cardinalidade do conjunto $X$, desejaríamos que obedecessem a:
\[
	\card{A} = \card{B} \text{ sse } \cardeq{A}{B}.
\]
Chamaremos a isto o Princípio \emph{de Cantor}, já que Cantor foi
provavelmente o primeiro a tê-lo claramente presente. (Diremos mais sobre a sua
relação com o Princípio de \emph{Hume} na \olref[hp]{sec}.) O nosso
objetivo é definir $\card{X}$, para cada $X$, de modo a que
satisfaça o Princípio de Cantor.

\end{document}
